\subsection{Questão 25}

\subsubsection{Enunciado}

Densidade de corrente de deslocamento uniforme. A Fig.~32-30 mostra uma região circular de raio 
\(R = 3{,}00\,\text{cm}\) na qual existe uma corrente de deslocamento dirigida para fora do papel.

A corrente de deslocamento possui uma densidade de corrente uniforme cujo valor absoluto é

\[
J_d = 6{,}00\,\text{A/m}^2.
\]

Determine o módulo do campo magnético produzido pela corrente de deslocamento

\begin{enumerate}
    \item[(a)] a \(2{,}00\,\text{cm}\) do centro da região;
    
    \item[(b)] a \(5{,}00\,\text{cm}\) do centro da região.
\end{enumerate}


\subsubsection{Resolução}

Pela lei de Ampere-Maxwell,

\[
B(2\pi r)=\mu_0 i_{d,\text{env}}.
\]

A densidade de corrente de deslocamento é uniforme, então a corrente de
deslocamento envolvida depende da área envolvida.

\begin{enumerate}
    \item[(a)] Para \(r=2{,}00\,\text{cm}<R\),

    \[
    i_{d,\text{env}}=J_d\pi r^2.
    \]

    Portanto,

    \[
    B(2\pi r)=\mu_0J_d\pi r^2,
    \]

    ou

    \[
    B=\frac{\mu_0J_dr}{2}.
    \]

    Assim,

    \[
    B=
    \frac{(4\pi\times 10^{-7})(6{,}00)(2{,}00\times 10^{-2})}{2}
    =
    7{,}54\times 10^{-8}\,\text{T}.
    \]

    \item[(b)] Para \(r=5{,}00\,\text{cm}>R\), a corrente envolvida é a
    corrente total da região:

    \[
    i_{d,\text{env}}=J_d\pi R^2.
    \]

    Logo,

    \[
    B=
    \frac{\mu_0J_dR^2}{2r}.
    \]

    Substituindo \(R=3{,}00\,\text{cm}\) e \(r=5{,}00\,\text{cm}\),

    \[
    B=
    \frac{(4\pi\times 10^{-7})(6{,}00)(3{,}00\times 10^{-2})^2}
    {2(5{,}00\times 10^{-2})}
    =
    6{,}79\times 10^{-8}\,\text{T}.
    \]
\end{enumerate}

\vspace{1cm}
